\chapter{Common-Sense Capture}
\label{ch:common-sense-capture}

\epigraph{The smart way to keep people passive and obedient is to strictly limit the spectrum of acceptable opinion, but allow very lively debate within that spectrum.}{Noam Chomsky}

\noindent
The most powerful doxonomic achievement is not persuading people to hold a belief.
It is making a belief so pervasive that it ceases to be recognized as a belief at all. A funded narrative becomes invisible infrastructure, indistinguishable from the texture of reality.

This chapter formalizes \emph{common-sense capture}\index{common-sense capture}: the phase transition by which a contested proposition becomes unquestioned background, the mechanisms that produce it, the conditions that sustain it, and the signatures by which it can be detected.


\section{Capture as Phase Transition}
\label{sec:ch8-phase-transition}

\begin{definition}[Common-Sense Capture]
\label{def:capture}
\index{common-sense capture!formal definition}
A belief $p$ is \emph{captured} in a population if:
\begin{enumerate}[nosep]
  \item $\mu(p) > \theta$ for a high threshold $\theta$ (widespread adoption),
  \item $\mathrm{Var}(\pi_j(p)) < \varepsilon$ for small $\varepsilon$ (tight concentration),
  \item the proposition $p$ is not salient as a \emph{claim}. It is experienced as a feature of reality.
\end{enumerate}
Conditions (1) and (2) are measurable.
Condition (3) is the distinctively doxonomic condition: $p$ has exited the space of contestation.
\end{definition}

The transition from ``widely held opinion'' to ``common sense'' is not gradual.
It exhibits threshold behavior analogous to phase transitions in physics.
Below the critical threshold, a belief is debated, contested, and recognized as a position.
Above it, the belief disappears into the background: it is no longer argued for because it no longer needs to be.

\begin{example}[Stages of Capture]
\label{ex:stages-capture}
Consider the proposition ``economic growth is the primary objective of government policy.''
\begin{enumerate}[nosep]
  \item \textbf{Contested (pre-1950s)}: alternative objectives were actively debated, including full employment, social solidarity, moral improvement, imperial expansion, and class harmony.
  \item \textbf{Dominant (1950s--1970s)}: GDP growth became the primary metric, but alternatives (quality of life, distributional equity) were still voiced in mainstream discourse.
  \item \textbf{Common sense (1980s--2000s)}: growth became so dominant that questioning it sounded eccentric. Policy debates assumed growth as the goal and argued only about means.
  \item \textbf{Captured/doxic (2000s--present)}: growth is not argued for; it is the unstated premise of virtually all policy discussion. A politician who says ``we should not pursue growth'' is not engaging in a legitimate policy debate; they are committing a category error.
\end{enumerate}
The transition from stage 2 to stage 4 is what doxonomics calls capture.
\end{example}


\section{Threshold Models}
\label{sec:ch8-thresholds}

Following \textcite{granovetter1978} and \textcite{schelling1978}, we model adoption as a threshold process.
The core insight is that individual adoption decisions depend on how many others have already adopted: people are more willing to hold a belief when they perceive it as socially supported.

\begin{model}[Granovetter-Schelling Belief Adoption]
\label{mod:granovetter}
\index{threshold model}
Each agent $j$ has a threshold $\tau_j \in [0,1]$ drawn from a distribution $G(\tau)$.
Agent $j$ adopts belief $p$ if and only if the fraction of the population already holding $p$ exceeds $\tau_j$:
\[
  \pi_j(p) = 1 \iff \mu(p) \geq \tau_j
\]
The equilibrium population fraction is a fixed point of
\[
  \mu^* = G(\mu^*)
\]
\end{model}

The threshold $\tau_j$ can be interpreted as agent $j$'s ``social proof requirement'': how many people must hold a belief before $j$ is willing to adopt it.
Agents with low thresholds are \emph{early adopters} (they adopt without much social support); agents with high thresholds are \emph{holdouts} (they adopt only when almost everyone else already has).

\begin{theorem}[Tipping Point]
\label{thm:tipping}
\index{tipping point}
If $G$ is continuous and $G(0) > 0$ (some agents adopt unconditionally), then there exists at least one stable equilibrium $\mu^*$.
A doxonomic investment that shifts $G$ downward (lowering thresholds by increasing social pressure, institutional endorsement, or reducing the perceived cost of adoption) can cause a discontinuous jump in $\mu^*$, a \emph{tipping point}.
\end{theorem}

\begin{proof}[Proof sketch]
The function $h(\mu) = G(\mu) - \mu$ is continuous with $h(0) = G(0) > 0$ and $h(1) = G(1) - 1 \leq 0$.
By the intermediate value theorem, $h$ has a zero.
Stability requires $G'(\mu^*) < 1$.
A downward shift in $G$ (uniform reduction in thresholds) raises $G(\mu)$ for all $\mu$, which can push $\mu^*$ past a saddle-node bifurcation, causing the low-adoption equilibrium to disappear and the system to jump to the high-adoption equilibrium.
\end{proof}

\begin{example}[Tipping in Belief Adoption]
\label{ex:tipping}
\index{tipping point!example}
Consider a population with threshold distribution $G(\tau) = \tau^{1/2}$ (many agents have low thresholds; few are extreme holdouts).
The fixed points satisfy $\mu = \mu^{1/2}$, giving $\mu^* = 0$ and $\mu^* = 1$.
The interior dynamics are: for any $\mu > 0$, we have $G(\mu) = \mu^{1/2} > \mu$, so adoption grows.
This system has only one stable equilibrium: $\mu^* = 1$ (complete adoption).
Any nonzero seed of adoption cascades to unanimity.

Now consider $G(\tau) = \tau^2$ (most agents have high thresholds).
The fixed points satisfy $\mu = \mu^2$, giving $\mu^* = 0$ and $\mu^* = 1$.
For $\mu \in (0,1)$, $G(\mu) = \mu^2 < \mu$, so adoption shrinks.
This system has only one stable equilibrium: $\mu^* = 0$ (no adoption).
Even substantial initial adoption erodes to zero.

The doxonomic implication: the \emph{shape} of the threshold distribution (determined by institutional environment, prior narrative exposure, and social structure) determines whether a belief can cascade.
Doxonomic investment works by reshaping $G$: lowering thresholds through institutional endorsement, repetition, and social proof until the system tips.
\end{example}

\begin{remark}[What Determines Thresholds?]
\label{rem:thresholds}
Individual thresholds $\tau_j$ are not fixed parameters.
They depend on:
\begin{itemize}[nosep]
  \item \textbf{Institutional endorsement}: when authoritative institutions adopt $p$, individual thresholds drop (``if Harvard says so, it must be reasonable'').
  \item \textbf{Perceived cost of dissent}: in environments where questioning $p$ is socially costly (mockery, professional risk, ostracism), thresholds drop (people adopt to avoid punishment, not because they are convinced).
  \item \textbf{Narrative coherence}: when $p$ fits into an existing ideological bundle (Definition~\ref{def:ideology}), thresholds drop because adoption is cognitively cheap.
  \item \textbf{Prior exposure}: repeated exposure to $p$ through multiple channels lowers thresholds via the illusory truth effect.
  \item \textbf{Identity alignment}: when $p$ is framed as aligned with a valued identity, thresholds drop; when framed as threatening, thresholds rise.
\end{itemize}
Doxonomic investment works on thresholds through all five mechanisms simultaneously.
\end{remark}


\section{The Doxonomic Investment Needed for Capture}
\label{sec:ch8-investment}

How much does capture cost?
This depends on the initial threshold distribution, the available channels, and the target belief.

\begin{model}[Capture Cost]
\label{mod:capture-cost}
\index{capture cost}
Let $G_0(\tau)$ be the initial threshold distribution and $G_F(\tau)$ the distribution after doxonomic investment $F$.
Assume investment uniformly lowers thresholds: $G_F(\tau) = G_0(\tau - \alpha F)$ for some effectiveness parameter $\alpha$.
The minimum investment $F^*$ needed for capture is the smallest $F$ such that the high-adoption fixed point of $\mu = G_F(\mu)$ exists and is globally stable.

At the tipping point, $G_F$ must be tangent to the 45-degree line from below:
\[
  G_F(\mu^*) = \mu^* \quad \text{and} \quad G_F'(\mu^*) = 1
\]
These two equations determine $(\mu^*, F^*)$.
\end{model}

\begin{remark}
The capture cost $F^*$ depends critically on the initial threshold distribution.
In populations with many low-threshold agents (already predisposed by prior narrative exposure or identity alignment), $F^*$ is small: a nudge suffices.
In populations with many high-threshold agents (resistant due to counter-institutions, lived experience, or strong alternative narratives), $F^*$ can be enormous.
This explains why the same doxonomic operation can succeed in one context and fail in another: the target population's threshold distribution matters as much as the budget.
\end{remark}


\section{Mechanisms of Capture}
\label{sec:ch8-mechanisms}

Capture does not happen through a single mechanism.
It typically involves the coordinated operation of several:

\subsection{Institutional Saturation}

When enough institutions treat $p$ as an operational premise, the cumulative effect is overwhelming.
A student who encounters $p$ in economics class, in management training, in the workplace, in advertising, in entertainment, and in political discourse receives $p$ from so many independent-seeming sources that it appears to be a natural fact rather than a specific claim.

\begin{definition}[Institutional Saturation]
\label{def:saturation}
\index{institutional saturation}
A belief $p$ is \emph{institutionally saturated} if the fraction of institutional contexts in which $p$ functions as an operational premise exceeds a threshold $\sigma$:
\[
  \frac{|\{U_i \in \mathcal{U} : p \in \sigma_i\}|}{|\mathcal{U}|} > \sigma
\]
where $\mathcal{U}$ is the cover of social life by institutional contexts and $\sigma_i$ is the local belief system of context $U_i$ (Chapter~\ref{ch:contradiction-coherence}).
\end{definition}

Institutional saturation is related to but distinct from population-level adoption.
A belief can be institutionally saturated ($p$ is embedded in workplaces, schools, and media) without being unanimously held (many individuals privately doubt $p$).
This is the pluralistic-ignorance configuration (Definition~\ref{def:pluralistic-ignorance}): the institutional environment enforces $p$ as operational reality even where private belief is weak.

\subsection{Crowding Out Alternatives}

Capture involves not just amplifying $p$ but \emph{suppressing alternatives}.
When the narrative infrastructure is saturated with $p$, there is no institutional space for $\neg p$: no funding, no platforms, no expert endorsement, no curriculum presence.
The alternatives do not need to be explicitly censored; they are simply not produced at sufficient volume to be encountered.

\begin{definition}[Narrative Crowding Out]
\label{def:crowding-out}
\index{crowding out}
A belief $p$ \emph{crowds out} alternative $q$ if the doxonomic investment in $p$ reduces the effective throughput for $q$:
\[
  \frac{\partial \Theta(q)}{\partial F(p)} < 0
\]
This occurs when $p$ and $q$ compete for the same infrastructure (media attention, curriculum space, expert endorsement, public attention) and $p$'s dominance leaves insufficient residual capacity for $q$.
\end{definition}

\subsection{Normalization Through Repetition}

Sheer repetition makes a claim feel true (the illusory truth effect).
When $p$ is encountered daily across multiple channels, it acquires the phenomenological quality of familiarity, which the cognitive system interprets as evidence for truth.
This is not a rational process; it is a cognitive heuristic that doxonomic systems exploit.

\subsection{Social Proof Cascades}

Once a critical mass of people (or, more importantly, a critical mass of \emph{visible} people) endorses $p$, social proof dynamics take over.
Each new adopter makes $p$ appear more normal, which lowers the threshold for the next adopter, which makes $p$ appear even more normal, a positive feedback cascade that accelerates toward capture.

\begin{model}[Social Proof Cascade]
\label{mod:cascade}
\index{social proof cascade}
Let $\mu(t)$ be the fraction holding $p$ at time $t$, with dynamics:
\[
  \frac{d\mu}{dt} = \alpha\phi(t)(1-\mu) + \beta\mu(1-\mu) - \delta\mu
\]
where $\alpha\phi(t)$ is institutional production, $\beta$ is the social proof coefficient, and $\delta$ is natural decay.
Setting $d\mu/dt = 0$, the non-trivial equilibria satisfy:
\[
  \mu^* = \frac{(\alpha\phi + \beta - \delta) \pm \sqrt{(\alpha\phi + \beta - \delta)^2 + 4\beta\alpha\phi}}{2\beta}
\]
When $\beta$ is large (strong social proof), the system is bistable: there is a low-$\mu$ equilibrium and a high-$\mu$ equilibrium, separated by an unstable threshold.
Doxonomic investment (increasing $\phi$) pushes the system across the threshold and into the high-$\mu$ basin.
\end{model}


\section{Path Dependence and Lock-In}
\label{sec:ch8-lock-in}

Once a belief crosses the tipping point, it becomes self-reinforcing through the feedback mechanisms described in Chapter~\ref{ch:belief-production-chain}: the belief generates behaviors that generate institutions that generate funding that generates more of the same belief.

\begin{definition}[Doxonomic Lock-In]
\label{def:lock-in}
\index{lock-in}
A belief $p$ is \emph{locked in} if the doxonomic equilibrium $\mu^*(p)$ is locally asymptotically stable and the basin of attraction is large, meaning that small perturbations (counter-narratives, contrary evidence, temporary funding cuts) do not dislodge it.
\end{definition}

\begin{proposition}[Lock-In Condition]
\label{prop:lock-in}
Lock-in occurs when the self-reinforcement gain exceeds the natural decay rate:
\[
  \frac{\partial \mu^{(t+1)}}{\partial \mu^{(t)}} \bigg|_{\mu = \mu^*} > 1 - \delta
\]
where $\delta$ is the rate at which belief decays absent reinforcement.
\end{proposition}

Lock-in has a temporal dimension.
The longer a belief has been locked in, the harder it is to dislodge, because:
\begin{itemize}[nosep]
  \item institutional infrastructure has been built around it (path dependence),
  \item professional careers depend on it (vested interests),
  \item curricula have transmitted it to new generations (generational reproduction),
  \item alternative beliefs have atrophied for lack of institutional support (crowding out),
  \item the belief has shaped behavior in ways that generate confirming evidence (self-fulfillment).
\end{itemize}

\begin{example}[Lock-In: Homo Economicus]
\label{ex:lock-in-homo-economicus}
The selfish-rational-agent model (homo economicus) has been locked in as the operational anthropology of economic policy since approximately the 1980s.
Dislodging it would require simultaneous intervention at multiple points:
\begin{itemize}[nosep]
  \item \textbf{Curricula}: rewriting economics textbooks to foreground conditional cooperation, behavioral complexity, and institutional context.
  \item \textbf{Professional norms}: changing what counts as ``rigorous'' economic analysis in journals, tenure review, and policy consulting.
  \item \textbf{Policy design}: replacing incentive-based regulation with trust-based institutional design.
  \item \textbf{Management theory}: replacing principal-agent models with cooperative organizational theories.
  \item \textbf{Popular culture}: producing narratives in which cooperative behavior is normal, not exceptional.
\end{itemize}
No single intervention suffices because the lock-in operates at all five levels simultaneously.
This is why lock-in is so durable: disrupting one level is compensated by the other four.
\end{example}


\section{Saturation and Diminishing Returns}
\label{sec:ch8-saturation-returns}

\begin{proposition}[Diminishing Returns Post-Capture]
\label{prop:diminishing}
Once common-sense capture is achieved ($\mu(p) > \theta$), the marginal belief change per additional dollar of funding approaches zero:
\[
  \lim_{\mu \to 1} \frac{\partial \mu}{\partial F} = 0
\]
Post-capture, doxonomic spending shifts from persuasion to \emph{maintenance}: preventing erosion rather than producing new converts.
\end{proposition}

This has several practical implications:

\begin{enumerate}[nosep]
  \item \textbf{Maintenance is cheaper than capture.}
  The cost of maintaining a captured belief is typically an order of magnitude lower than the cost of capturing it.
  Once the infrastructure is built, the institutions staffed, and the curricula written, the system runs on inertia plus modest annual investment.
  This asymmetry explains why dominant narratives persist long after the original funding rationale has changed.

  \item \textbf{Over-investment is common.}
  Funders often continue spending at capture-level rates even after capture has been achieved, because they cannot easily measure the current belief level and err on the side of caution.
  This produces \emph{doxonomic waste}: resources spent reinforcing beliefs that are already unquestioned.

  \item \textbf{The marginal dollar is better spent on new frontiers.}
  Once $p$ is captured, a rational doxonomic investor should redirect resources to adjacent beliefs ($p_2, p_3, \ldots$) that extend the ideological bundle, rather than pouring more money into $p$ itself.
  This is the doxonomic analogue of market saturation in product economics.
\end{enumerate}


\section{Detecting Capture: Empirical Signatures}
\label{sec:ch8-detection}

How does a researcher identify that a belief has been captured rather than merely being popular?
The detection methods from Section~\ref{sec:ch2-detecting-doxa} apply, but we can add quantitative signatures:

\begin{definition}[Capture Signature]
\label{def:capture-signature}
\index{capture!detection}
A belief $p$ shows a \emph{capture signature} if:
\begin{enumerate}[nosep]
  \item \textbf{Asymmetric argumentation}: $p$ appears overwhelmingly as a premise rather than a conclusion in public discourse (measurable via corpus analysis).
  \item \textbf{Low narrative entropy}: the narrative space around $p$'s domain has collapsed to a small number of frames ($H(t)$ below baseline; see Definition~\ref{def:narrative-entropy}).
  \item \textbf{High institutional saturation}: $p$ functions as an operational assumption in a large fraction of institutional contexts (Definition~\ref{def:saturation}).
  \item \textbf{Anomaly response}: challenges to $p$ are met with confusion or ridicule rather than counter-argument.
  \item \textbf{Funding-belief decoupling}: $p$'s prevalence persists even during periods of reduced doxonomic investment (the belief runs on institutional inertia, not active funding).
  \item \textbf{Cross-domain consistency}: $p$ appears in economics, management, law, media, and everyday speech without any domain citing the others as its source.
\end{enumerate}
\end{definition}

\begin{remark}
Not all captured beliefs are false, and not all uncaptured beliefs are true.
The capture signature identifies a \emph{structural} property of the belief's social position, not an epistemic property of its content.
``The earth orbits the sun'' is also captured (near-universal, low variance, invisible as a claim), but it is captured \emph{because it is true and well-evidenced}, not because of doxonomic investment.
The researcher must distinguish evidence-driven capture from investment-driven capture, which requires the separation of audit and evaluation (Principle~\ref{prin:separation}).
\end{remark}


\section{Historical Case: The Capture of ``Markets Know Best''}
\label{sec:ch8-case}

The belief that markets are generally superior to government in allocating resources provides a detailed case study of the capture process.

\textbf{Phase 1: Contestation (1930s--1960s).}
In the aftermath of the Great Depression, market-skeptical views dominated: Keynesian macroeconomics, strong regulation, public ownership of utilities, and active industrial policy were mainstream.
The belief ``markets know best'' was a minority position, associated with a small network of Austrian and Chicago economists.

\textbf{Phase 2: Infrastructure Building (1947--1980).}
The Mont P\`{e}lerin Society (1947), the Foundation for Economic Education, and later the Heritage Foundation (1973), Cato Institute (1977), and dozens of smaller think tanks built the narrative infrastructure.
Funding from Olin, Scaife, Koch, and Bradley foundations supported law-and-economics programs at universities, endowed chairs, and produced a generation of economists, lawyers, and policy intellectuals committed to market-first thinking \autocite{mirowski2009, slobodian2018}.

\textbf{Phase 3: Political Capture (1979--1990).}
Thatcher and Reagan brought the narrative to political power.
Deregulation, privatization, and tax cuts were enacted.
The behavioral consequences (rising stock markets, new consumer products, visible wealth creation for some) appeared to confirm the narrative.

\textbf{Phase 4: Common-Sense Capture (1990s--2008).}
By the 1990s, even center-left parties (Clinton's Democrats, Blair's New Labour) adopted market-first assumptions.
The ``Washington Consensus'' made market liberalization the default prescription for developing countries.
Alternatives (industrial policy, public ownership, cooperative models) were no longer debated; they were simply absent from serious policy discussion.

\textbf{Phase 5: Partial Fracture (2008--present).}
The financial crisis disrupted the capture by generating a massive contradiction (markets spectacularly failed to self-correct).
But the fracture was partial: despite the crisis, the policy response (austerity, quantitative easing for banks, limited re-regulation) largely preserved market-first assumptions.
The belief was shaken but not dislodged: evidence of deep lock-in.

\begin{remark}
This case illustrates the full capture trajectory: from minority position to infrastructure building to political implementation to common-sense invisibility to partial fracture under contradiction.
The total doxonomic investment over 60 years was plausibly in the billions of dollars.
The capture was not permanent (the 2008 crisis opened cracks), but the lock-in was strong enough to survive even a dramatic empirical refutation.
\end{remark}


\section{Escape from Capture}
\label{sec:ch8-escape}

If common-sense capture is so durable, how does it ever end?
Chapter~\ref{ch:how-belief-regimes-fracture} treats this in detail, but a preview is useful here.

\begin{conjecture}[Conditions for De-Capture]
\label{conj:de-capture}
\index{de-capture}
A captured belief $p$ can be dislodged when:
\begin{enumerate}[nosep]
  \item \textbf{Contradiction accumulates}: the gap between $p$ and lived experience becomes too large for institutional maintenance to bridge (e.g., ``markets self-correct'' during a prolonged recession).
  \item \textbf{Counter-infrastructure exists}: alternative institutions have been built that can supply a replacement narrative at sufficient volume (e.g., the environmental movement built enough infrastructure to challenge growth-as-progress).
  \item \textbf{Generational turnover}: new cohorts arrive without the formative exposure that locked in $p$ for the previous generation.
  \item \textbf{Elite defection}: high-credibility insiders publicly break with $p$, reducing the social cost of dissent (e.g., prominent economists questioning homo economicus after 2008).
  \item \textbf{Crisis opens the Overton window}: an exogenous shock temporarily suspends the normal inertia of common sense (Chapter~\ref{ch:crisis-doxonomics}).
\end{enumerate}
Typically, multiple conditions must coincide.
Single-condition de-capture is rare; this is why lock-in is so durable.
\end{conjecture}


\section{Notes and References}
\label{sec:ch8-notes}

Threshold models originate with \textcite{granovetter1978} and \textcite{schelling1978}.
Cascade and herding models appear in \textcite{banerjee1992} and \textcite{bikhchandani1992}.
Path dependence in institutional contexts is treated in \textcite{north1990}.
The concept of lock-in from technology economics \autocite{arthur1989} transfers to the belief domain.
On the invisibility of dominant frames, see \textcite{lukes2005}.

The capture of market-first ideology is documented in \textcite{mirowski2009}, \textcite{harvey2005}, \textcite{slobodian2018}, and \textcite{phillipsfein2009}.
On the illusory truth effect and repetition-based persuasion, see \textcite{kahneman2011}.
Social proof dynamics are analyzed in \textcite{cialdini2006}.
On the partial fracture of market fundamentalism after 2008, see \textcite{streeck2016}.
For the concept of crowding out in media and attention markets, see \textcite{mcchesney1999}.
On narrative entropy as a measure of discursive openness, see the information-theoretic framework in Chapter~\ref{ch:text-media}.
The concept of institutional saturation extends the Althusserian notion of ideological state apparatuses \autocite{althusser1970} into a measurable condition.


\section{Exercises}

\begin{exercise}
Given a threshold distribution $G(\tau) = \tau^2$ on $[0,1]$, find all fixed points of $\mu^* = G(\mu^*)$.
Determine their stability.
What is the minimum uniform threshold reduction $\Delta$ needed to create a stable high-adoption equilibrium?
Plot $G(\tau)$ and the shifted $G(\tau - \Delta)$ to illustrate.
\end{exercise}

\begin{exercise}
For the social proof cascade model (Model~\ref{mod:cascade}) with $\alpha = 0.3$, $\beta = 0.5$, $\delta = 0.1$, find the minimum institutional investment $\phi^*$ needed for a stable high-adoption equilibrium to exist.
Plot the phase portrait for $\phi$ just below and just above $\phi^*$.
\end{exercise}

\begin{exercise}
Estimate the maintenance cost of the belief ``economic growth is always desirable'' by identifying three institutions that continuously reinforce it and estimating their annual budgets.
Compare the maintenance cost to a rough estimate of the original capture cost (infrastructure built from the 1940s onward).
\end{exercise}

\begin{exercise}
Design a formal test to distinguish a belief that is common sense (Definition~\ref{def:common-sense}) from one that is captured (Definition~\ref{def:capture}).
What empirical evidence would differentiate the two?
Apply your test to ``democracy is the best form of government.''
\end{exercise}

\begin{exercise}
\label{ex:ch8-crowding}
Model the crowding-out effect formally.
Let two beliefs $p$ and $q$ compete for a shared narrative infrastructure with total capacity $\Theta_{\max}$.
Show that if $\Theta(p) + \Theta(q) \leq \Theta_{\max}$, increased investment in $p$ necessarily reduces throughput for $q$.
Under what conditions can $q$ survive despite being outspent?
\end{exercise}

\begin{exercise}
Apply the capture signature (Definition~\ref{def:capture-signature}) to the belief ``a university degree is necessary for a good career.''
Assess each of the six criteria.
Is this belief captured?
If so, who benefits from its capture?
\end{exercise}

\begin{exercise}
\label{ex:ch8-de-capture}
Choose a belief that has been partially de-captured in the past 20 years (e.g., ``smoking is glamorous,'' ``homosexuality is a disorder,'' ``colonialism was civilizing'').
Trace the de-capture process through the five conditions in Conjecture~\ref{conj:de-capture}.
Which conditions were most important?
Was the de-capture complete or partial?
\end{exercise}

\begin{exercise}
The chapter claims that ``maintenance is cheaper than capture.''
Formalize this claim using the social proof cascade model.
Show that the per-period funding needed to maintain $\mu^*$ near 1 is lower than the total funding needed to reach $\mu^*$ from $\mu = 0$.
Under what conditions is this \emph{not} true?
\end{exercise}
